\section*{Chapter \ref{ch:b1:optical-mineralogy} --- Optical Mineralogy}

\begin{solution}{exo:b1:optical-mineralogy:1}
Quartz: $\qty{30000}{nm}\times0.009 = \qty{270}{nm}$, a \omterm{def:b1:optical-mineralogy:interference-colour}{first-order} white.
Calcite: $\qty{30000}{nm}\times0.172 = \qty{5160}{nm}$, far beyond the chart:
a pale \omterm{def:b1:optical-mineralogy:interference-colour}{high-order} white, faintly pink or green.
\end{solution}

\begin{solution}{exo:b1:optical-mineralogy:2}
$I/I_{\max} = \sin^2 2\theta$: $0$ at \ang{0}; $\sin^2\ang{45} = 0.5$ at
\ang{22.5}; $1$ at \ang{45}; $0.5$ at \ang{67.5}; $0$ at \ang{90}.
\end{solution}

\begin{solution}{exo:b1:optical-mineralogy:3}
\omterm{def:b1:rock-forming-minerals:groups}{Garnet} is cubic, hence isotropic: light leaves it vibrating as it entered,
along the polariser, and the analyser stops it; $\Delta = 0$ in every
position. A quartz grain cut perpendicular to its \omterm{def:b1:optical-mineralogy:optic-axis}{optic axis} also has
$n_1 = n_2$ for light travelling along that axis, so $\Delta = 0$: it too stays
dark. In \omterm{def:b1:field-lab-methods-1:microscope}{plane-polarised light} the \omterm{def:b1:rock-forming-minerals:groups}{garnet} has very high \omterm{def:b1:optical-mineralogy:relief}{relief} ($1.83$
against $1.54$), the quartz almost none.
\end{solution}

\begin{solution}{exo:b1:optical-mineralogy:4}
$|n - 1.54|$: almandine $0.29$, forsterite $0.11$, orthoclase $0.02$, quartz
$0.01$; \omterm{def:b1:optical-mineralogy:relief}{relief} in that order, from highest to lowest. Orthoclase has the lower
index: negative \omterm{def:b1:optical-mineralogy:relief}{relief}, and its \omterm{def:b1:optical-mineralogy:relief}{Becke line} moves out of the grain, into the
mount, when the focus is raised.
\end{solution}

\begin{solution}{exo:b1:optical-mineralogy:5}
$\delta = 650/30000 = 0.022$. A green \omterm{def:b1:optical-mineralogy:pleochroism}{pleochroic} \omterm{def:b1:rock-forming-minerals:mineral}{mineral} with cleavages at
\ang{56} and \ang{124} and a \omterm{def:b1:optical-mineralogy:birefringence}{birefringence} near $0.020$: hornblende.
\end{solution}

\begin{solution}{exo:b1:optical-mineralogy:6}
$\tan\ang{30} = 0.5774$. Quartz: $(0.013/1.54)\times0.5774 = \qty{0.00487}{rad}
= \ang{0.28}$. Forsterite: $(0.111/1.54)\times0.5774 = \qty{0.0416}{rad} =
\ang{2.38}$. Ratio $0.111/0.013 = 8.5$. Exactly: $\sin r = 1.54\sin\ang{30}/1.651
= 0.4664$, $r = \ang{27.80}$, a deviation of \ang{2.20}: the first-order
formula overestimates by \qty{8}{\percent} for so large an index difference.
\end{solution}

\begin{solution}{exo:b1:optical-mineralogy:7}
$270 + 540 = \qty{810}{nm}$, a \omterm{def:b1:optical-mineralogy:interference-colour}{second-order} green, when the slow directions of
the grain and the plate are parallel (the delays add); $|270 - 540| =
\qty{270}{nm}$, \omterm{def:b1:optical-mineralogy:interference-colour}{first-order} white again, when the slow direction of the grain
is parallel to the fast direction of the plate.
\end{solution}

\begin{solution}{exo:b1:optical-mineralogy:8}
At full thickness $\Delta = \qty{30000}{nm}\times0.052 = \qty{1560}{nm}$. From
the edge ($\Delta = 0$) inwards, the \omterm{def:b1:optical-mineralogy:retardation}{retardation} passes 542 and
\qty{1074}{nm} but not \qty{1605}{nm}: two red-violet fringes.
\end{solution}

\begin{solution}{exo:b1:optical-mineralogy:9}
The \omterm{def:b1:optical-mineralogy:extinction}{extinction angle} measured on a grain depends on how the grain is cut:
only sections in a particular orientation show the largest angle, and others
show less, down to zero. The largest value, about \ang{37}, is the one to
compare: it matches diopside (about \ang{38}), not hornblende (at most
\ang{19}).
\end{solution}

\begin{solution}{exo:b1:optical-mineralogy:10}
With the analyser along $x$, the $x$ components of the two waves are
$E_0\cos\theta\cdot\cos\theta$ and $-E_0\sin\theta\,e^{-i\varphi}\cdot(-\sin\theta)$,
so $A = E_0(\cos^2\theta + \sin^2\theta\,e^{-i\varphi})$ and
\[
  |A|^2 = E_0^2\bigl(\cos^4\theta + \sin^4\theta + 2\sin^2\theta\cos^2\theta\cos\varphi\bigr)
  = E_0^2\bigl(1 - 2\sin^2\theta\cos^2\theta\,(1 - \cos\varphi)\bigr),
\]
using $\cos^4\theta + \sin^4\theta = 1 - 2\sin^2\theta\cos^2\theta$. Since
$2\sin^2\theta\cos^2\theta = \frac12\sin^2 2\theta$ and $1 - \cos\varphi =
2\sin^2(\varphi/2)$, $I_\parallel = I_0(1 - \sin^2 2\theta\sin^2(\pi\Delta/\lambda))$.
At each wavelength, $I_\parallel + I_\perp = I_0$: the colours between parallel
polars are the complements of those between \omterm{def:b1:field-lab-methods-1:microscope}{crossed polars}. A grain of
\qty{550}{nm}, magenta between \omterm{def:b1:field-lab-methods-1:microscope}{crossed polars}, is pale yellowish green
between parallel polars.
\end{solution}

\begin{solution}{exo:b1:optical-mineralogy:11}
(a) $\Delta/\lambda$ runs from $5000/700 = 7.14$ to $5000/400 = 12.5$; it
passes the whole numbers 8 to 12: five zeros. (b) The zeros are at
$\lambda_k = \Delta/k$; two successive ones are $\Delta/k - \Delta/(k + 1) =
\Delta/(k(k + 1)) \approx \lambda^2/\Delta$ apart, that is $550^2/5000 =
\qty{60.5}{nm}$ near \qty{550}{nm}. (c) Five narrow dark bands spread across
the spectrum remove a little of every colour, and the rest passes with an
average weight of one half: the eye sees a pale, nearly neutral tint.
\end{solution}

\begin{solution}{exo:b1:optical-mineralogy:12}
The mean of $\sin^2 2\theta$ over a turn is $\frac1{2\pi}\int_0^{2\pi}\sin^2 2\theta\,
\mathrm d\theta = \frac12$, so the mean intensity is
$\frac12 I_0\sin^2(\pi\Delta/\lambda)$. She computes $\delta = \Delta/30000$
instead of $\Delta/25000$ (in nm): all her values are $25/30 = 0.83$ of the
true ones, too low by \qty{17}{\percent}.
\end{solution}

\begin{solution}{pb:b1:optical-mineralogy:1}
\textbf{1.} B: dark rims mean high \omterm{def:b1:optical-mineralogy:relief}{relief}; A's barely visible boundaries mean
low \omterm{def:b1:optical-mineralogy:relief}{relief}.

\textbf{2.} A: $(1.544 + 1.553)/2 = 1.5485$, $|n - n_m| = 0.0085$. B:
$(1.664 + 1.672 + 1.694)/3 = 1.6767$, $|n - n_m| = 0.1367$.

\textbf{3.} The deviations are proportional to $|n - n_m|$: $0.1367/0.0085 =
16.1$.

\textbf{4.} B has a higher index than the mount: positive \omterm{def:b1:optical-mineralogy:relief}{relief}.

\textbf{5.} \omterm{def:b1:optical-mineralogy:pleochroism}{Pleochroism}; biotite.

\textbf{6.} Two cleavages at \ang{56} and \ang{124}: an \omterm{def:b1:rock-forming-minerals:groups}{amphibole}.

\textbf{7.} $\qty{30000}{nm}\times0.009 = \qty{270}{nm}$: \omterm{def:b1:optical-mineralogy:interference-colour}{first-order} white.

\textbf{8.} Each grain is cut in its own orientation and shows a value of
$|n_1 - n_2|$ between 0 and $\delta$; only grains cut parallel to the \omterm{def:b1:optical-mineralogy:optic-axis}{optic
axis} show $\delta$ itself. The highest colour among many grains is the one to
use.

\textbf{9.} $d = 342/0.009 = \qty{38000}{nm} = \qty{38}{\micro m}$.

\textbf{10.} $(38 - 30)/30 = \qty{26.7}{\percent}$.

\textbf{11.} $\sin^2(\pi\times342/550) = 0.861$; at \qty{30}{\micro m},
$\sin^2(\pi\times270/550) = 0.999$.

\textbf{12.} $\qty{38000}{nm}\times0.008 = \qty{304}{nm}$ to
$\qty{38000}{nm}\times0.013 = \qty{494}{nm}$: white to orange-red, all below
the end of the first \omterm{def:b1:optical-mineralogy:interference-colour}{order} at \qty{542}{nm}.

\textbf{13.} Between 1074 and \qty{1605}{nm}: third \omterm{def:b1:optical-mineralogy:interference-colour}{order}.

\textbf{14.} $\delta = 1140/38000 = 0.030$.

\textbf{15.} $1140/30000 = 0.038$, which would suggest an \omterm{def:b1:rock-forming-minerals:groups}{olivine}. A
magnesian \omterm{def:b1:rock-forming-minerals:groups}{olivine} and quartz react to give \omterm{def:b1:rock-forming-minerals:groups}{pyroxene},
\ce{Mg2SiO4 + SiO2 -> 2 MgSiO3}; a rock that has crystallised with quartz
should not keep both.

\textbf{16.} $0.030$ is the \omterm{def:b1:optical-mineralogy:birefringence}{birefringence} of diopside; its high \omterm{def:b1:optical-mineralogy:relief}{relief}, its
lack of colour and its two cleavages near \ang{87} agree: B is a
clinopyroxene, diopside.

\textbf{17.} $\qty{30000}{nm}\times0.030 = \qty{900}{nm}$: \omterm{def:b1:optical-mineralogy:interference-colour}{second-order} yellow.

\textbf{18.} By \qty{8}{\micro m}; quartz then shows at most \qty{270}{nm},
\omterm{def:b1:optical-mineralogy:interference-colour}{first-order} white.

\textbf{19.} \omterm{def:b1:optical-mineralogy:twin}{Polysynthetic twinning}: D is \omterm{def:b1:rock-forming-minerals:groups}{plagioclase}.

\textbf{20.} The composition of the \omterm{def:b1:rock-forming-minerals:groups}{plagioclase} (its anorthite content).

\textbf{21.} An \omterm{def:b1:rock-forming-minerals:groups}{amphibole} with an \omterm{def:b1:optical-mineralogy:extinction}{extinction angle} of \ang{17}: hornblende
(\ang{15}--\ang{19}), not a \omterm{def:b1:rock-forming-minerals:groups}{pyroxene} (\ang{38} for diopside).

\textbf{22.} $\qty{38000}{nm}\times0.020 = \qty{760}{nm}$: a \omterm{def:b1:optical-mineralogy:interference-colour}{second-order}
green.

\textbf{23.} It is cut perpendicular to an \omterm{def:b1:optical-mineralogy:optic-axis}{optic axis}: its two indices are
equal for light travelling along the axis, $\Delta = 0$, and it stays dark like
an isotropic grain.

\textbf{24.} A quartz, B diopside, C biotite, D \omterm{def:b1:rock-forming-minerals:groups}{plagioclase}, E hornblende: an
intermediate igneous rock with some quartz, such as a quartz diorite.

\textbf{25.} The section is \textbf{\qty[text-series-to-math=true]{38}{\micro m}} thick.
\end{solution}
